%% notes-tex/database/kg-build-system/sections/1-overview.tex

\section{The system in one figure\secstatus{todo}}
\label{sec:overview}

%% SOURCE: notes/kg-build-system.mermaid.build-flow.md, rendered by
%% notes/assets/publish/scripts/mermaid_pdf.sh and scaled to the text block by
%% `make diagrams`; the boxes name the scripts under scripts/ and database/slurm/scripts/.
\tcfig{figures/build-flow.pdf}{From source to release artifact. Yellow boxes are
files in git or the development data tree; orange boxes are steps that run on the
build host; the purple box is the release artifact; the blue box is the store the
build host serves; red hexagons are gates that refuse rather than degrade. The
cluster is the one box that moves when the build moves.}{fig:flow}

%% SOURCE: notes/kg-build-system.mermaid.serve-flow.md, rendered by
%% notes/assets/publish/scripts/mermaid_pdf.sh and scaled to the text block by
%% `make diagrams`; the steps are those of scripts/kg_release.sh and scripts/ops.sh.
\tcfig{figures/serve-flow.pdf}{From release artifact to a client. The archive on
Taiga holds every release as verified files; \texttt{deploy} restores one into any
serving host; \texttt{ops.sh sync} fails while hosts differ; the client gate
refuses a store the installed package does not fully reproduce. Same palette as
Fig.~\ref{fig:flow}.}{fig:serve}

\notesec{Five stages}
A release passes through five stages, each with one owner and one output
(Table~\ref{tab:stages}). The source stage lives in git on every machine. The
build stage is host-bound: it needs the development LMDB tree, docker, slurm and a
fast local disk, and today that host is GilaHyper. The archive stage is Taiga,
where a release is a directory of files and nothing runs. The serving stage is any
host with a Neo4j container on local disk, and the client stage is whatever
process opens a bolt connection with \texttt{torchcell} installed.

%% SOURCE: hand-authored from scripts/kg_release.sh, scripts/ops.sh,
%% database/slurm/scripts/gilahyper_live_rebuild-slurm_docker.slurm and
%% torchcell/knowledge_graphs/releases.py.
\begin{table}[htbp]
\centering
\footnotesize
\caption{The five stages, their owner, and the gate each one enforces.}
\label{tab:stages}
\begin{tabular}{l>{\raggedright\arraybackslash}p{0.22\textwidth}>{\raggedright\arraybackslash}p{0.26\textwidth}>{\raggedright\arraybackslash}p{0.26\textwidth}}
\toprule
Stage & Runs on & Produces & Gate \\
\midrule
Source & git, every machine & a tagged commit \texttt{vX.Y.Z}; the committed snapshot \texttt{database/releases/<release>.json} & semantic-release bumps only on \texttt{REL}, \texttt{DB}, \texttt{PATCH}, \texttt{API} \\
Build & the build host, under slurm & the store on \texttt{/db}, the \texttt{KgRelease} node, the manifest, the snapshot written into the checkout & the build commit must carry a package tag; every dev LMDB must be fresh \\
Archive & the build host, then Taiga & \texttt{<release>/} with the \texttt{.backup}, \texttt{kg\_manifest.json}, \texttt{release.json}, \texttt{SHA256SUMS} & \texttt{sha256sum -c} on arrival \\
Serve & any host with a Neo4j container on local disk & database \texttt{kg-<release>}, aliases \texttt{latest} and \texttt{pinned} & the restored node must match \texttt{release.json}; \texttt{ops.sh sync} fails on divergence \\
Client & any process with \texttt{torchcell} installed & an LMDB of query records & \texttt{require\_paired}: every served dataset's closure reproduced, else refused \\
\bottomrule
\end{tabular}
\end{table}

\notesec{What is host-independent}
Everything in Fig.~\ref{fig:serve}. The backup file
is Neo4j's own online-backup format, the archive directory is plain files with
checksums, and \texttt{kg\_release.sh deploy} runs on whichever host has the
container, reading the archive root from an environment variable. Moving the
build changes the second column of Table~\ref{tab:stages} for one row.
